% !TEX root = ../main.tex
\clearpage
\phantomsection
\section*{10 \textbf{Dissertation/Thesis Outline}}
\label{sec:dissertationthesis-outline}
\addcontentsline{toc}{section}{10 Dissertation/Thesis Outline}

The Dissertation consists of the following six chapters:


\begin{itemize}
\item[\textendash] Chapter 1: Introduction and research problem. This chapter will refine the proposal into the thesis introduction by presenting the DeFi lending context, the on-chain wallet reputation problem, the research questions, objectives, scope, significance, and limitations.
\item[\textendash] Chapter 2: Literature review and theoretical framework. This chapter will expand the proposal literature review into a comprehensive analysis of graph-based reputation scoring, DeFi credit-risk prediction, ERC-20 allowance semantics, PageRank variants, and theoretical foundations for reputation structure, predictive validity, and computational efficiency.
\item[\textendash] Chapter 3: Research methodology. This chapter will specify the philosophical position, research design, data sources, sampling logic, feature engineering, graph-construction procedures, baseline reproduction, statistical tests, robustness checks, ethics clearance, and reproducibility controls.
\item[\textendash] Chapter 4: Implementation and empirical results. This chapter will report on the data pipeline, EndorseRank implementation, baseline implementations, dataset characteristics, graph statistics, predictive-model results, rank-divergence findings, and computational benchmark outcomes.
\item[\textendash] Chapter 5: Discussion. This chapter will interpret the findings against each research question and objective, explain the theoretical and practical implications for DeFi risk assessment, evaluate threats to validity, and discuss how the results support or limit the claimed contribution.
\item[\textendash] Chapter 6: Conclusion and future work. This chapter will synthesize the study, state the original contribution to knowledge, revisit the research questions, summarize limitations, and propose future work on cross-protocol validation, adversarial manipulation, and broader on-chain trust signals.
\end{itemize}
\phantomsection
